% paper_review.tex (NeurIPS-style review draft; single-file; no TOC)
\documentclass{article}

% NeurIPS style file (place neurips_2025.sty in the same folder)
%\usepackage{neurips_2025}
\usepackage[margin=1in]{geometry}

\usepackage{amsmath,amssymb}
\usepackage{booktabs}
\usepackage[hypertexnames=false]{hyperref}
\usepackage{microtype}
\usepackage{graphicx}
\usepackage{enumitem}
\usepackage{algorithm}
\usepackage{algorithmic}

\title{Collapse-Resilient Decoding via Lightweight Online Intervention:%
\texorpdfstring{\\}{ }A Paired-Seed Study Across Five Prompt Domains}

% Review (anonymous). Camera-ready: replace with your name.
\author{Anonymous authors}

\begin{document}
\maketitle

\begin{abstract}
Long-form generation with locally-hosted language models can suffer from \emph{collapse},
including degenerative repetition and premature termination.
We present a lightweight decoding-time intervention scheme that monitors online signals
during generation and performs minimal corrective actions when collapse is detected.
We evaluate on Gemma 7B (Ollama) across five prompt domains (HUM/STEM/ETH/TEMP/META)
using a paired-seed design (500 seeds per pillar, 512 tokens fixed).
Across all pillars, controlled decoding substantially reduces collapse with small runtime overhead,
and all summary artifacts are produced via an auditable pipeline from raw JSONL logs to Table~1.
\end{abstract}

\section{Introduction}
Open and locally-hosted LLMs are increasingly used for long-form generation,
yet decoding can exhibit \emph{collapse} such as repetition loops or early termination.
Rather than modifying model weights or retraining, we focus on inference-time robustness:
monitor signals online and intervene only when collapse becomes likely.

\paragraph{Contributions.}
\begin{enumerate}[leftmargin=*, itemsep=2pt]
\item A minimal, model-agnostic online monitoring + intervention method operating purely at inference time.
\item A strict paired-seed evaluation protocol across five prompt domains with auditable artifacts.
\item Empirical evidence of large collapse-rate reductions with modest runtime overhead (Table~\ref{tab:main}).
\end{enumerate}

\section{Method}
\subsection{Domains (pillars) and paired-seed protocol}
We define five prompt domains (pillars): HUM, STEM, ETH, TEMP, and META.
For each pillar, we run seeds $s \in \{0,\dots,499\}$ in two modes: \textbf{baseline} and \textbf{controlled}.
Paired seeds ensure both modes use the same seed and deterministic prompt selection, enabling paired inference.

\subsection{Online collapse detectors}
Generation is monitored in segments. Each segment produces signals used for detection:
\begin{itemize}[leftmargin=*]
  \item \textbf{$\pi$-flag (calibrated scalar risk score):}
  $\pi$ is \emph{not} a probability; it is a calibrated scalar proxy for degeneration risk computed per segment,
  normalized using fixed calibration quantiles (e.g., $q_{50}, q_{95}$) and clamped to $[0,1]$.
  A collapse trajectory is flagged if $\pi > \epsilon_{\pi}$ for $k$ consecutive monitoring windows.
  \item \textbf{REP-flag:} an $n$-gram repetition detector; if repetition score exceeds a threshold, REP triggers.
  \item \textbf{TOO\_SHORT:} if total generated tokens fall below \texttt{min\_tokens}, the run is invalid/short.
\end{itemize}
The run-level collapse indicator is the logical OR of these conditions, using the same rule set for both modes.

\subsection{Intervention policy (explicit)}
Baseline differs from controlled \emph{only} by whether interventions are allowed.
\textbf{Baseline:} \texttt{max\_interventions=0}. \;
\textbf{Controlled:} allow up to \texttt{max\_interventions=100}.
When a collapse trajectory is detected, we apply a minimal corrective action that preserves the same prompt/seed:
we \emph{restart the current segment} (re-decode that segment) with a fixed ``repair'' decoding tweak
(temperature downshift and/or top-$p$ downshift for that segment only), then resume normal settings.
All intervention events are logged.

\begin{algorithm}[t]
\caption{Online monitoring and lightweight intervention (controlled mode)}
\label{alg:policy}
\begin{algorithmic}[1]
\STATE \textbf{Inputs:} prompt, seed $s$, max tokens $T$, segment length $L$, thresholds $(\epsilon_\pi, k)$, REP params, budget $B$
\STATE Initialize collapse\_streak $\leftarrow 0$, interventions $\leftarrow 0$, tokens\_generated $\leftarrow 0$
\WHILE{tokens\_generated $< T$}
  \STATE Decode next segment of up to $L$ tokens under base settings $\theta$ (temp/top-$p$ fixed)
  \STATE Compute signals for this segment: $\pi$, REP, and update tokens\_generated
  \IF{$\pi > \epsilon_\pi$} \STATE collapse\_streak $\leftarrow$ collapse\_streak $+ 1$ \ELSE \STATE collapse\_streak $\leftarrow 0$ \ENDIF
  \STATE collapse\_flag $\leftarrow$ (collapse\_streak $\ge k$) OR (REP triggers)
  \IF{collapse\_flag AND interventions $< B$}
     \STATE interventions $\leftarrow$ interventions $+ 1$
     \STATE \textbf{Intervene:} re-decode \emph{the current segment} under repair settings $\theta'$ (e.g., lower temp/top-$p$),
            then resume base $\theta$
     \STATE collapse\_streak $\leftarrow 0$
  \ENDIF
\ENDWHILE
\STATE \textbf{Finalize:} TOO\_SHORT if tokens\_generated $< \texttt{min\_tokens}$; log runtime and interventions
\end{algorithmic}
\end{algorithm}

\subsection{Calibration and ``Sensitive'' setting}
The $\pi$ detector uses a calibration file specific to the Gemma 7B (Ollama) setup.
We use a ``Sensitive'' calibration artifact (not publicly distributed as raw text) while keeping the full pipeline,
hashes/paths, and derived summary outputs auditable.

\section{Experimental Setup}
\subsection{Model and decoding}
Model: \texttt{gemma:7b} via Ollama.
Decoding is fixed across all runs: \texttt{max\_new\_tokens=512}, \texttt{temperature=1.2}, \texttt{top\_p=0.95}, \texttt{max\_segments=32}.
Detection hyperparameters are fixed (e.g., $\epsilon_\pi=0.85$, $k=3$, repetition settings).

\subsection{Primary metrics and statistics}
Primary outcome is \textbf{collapse rate} (fraction of runs flagged as collapsed).
We report absolute difference $\Delta = p_{\mathrm{ctrl}} - p_{\mathrm{base}}$ and effect sizes: risk ratio (RR) and odds ratio (OR).
For confidence intervals:
\begin{itemize}[leftmargin=*]
  \item Collapse rates use 95\% Wilson intervals.
  \item $\Delta$/RR/OR use paired bootstrap over seeds (B=20{,}000).
  \item RR/OR apply Haldane--Anscombe correction (+0.5) when any cell count is zero.
\end{itemize}

\subsection{Lightweight baseline comparator (decode-only)}
To address whether simple decode-time heuristics suffice, we include a minimal comparator:
\textbf{no-repeat 3-gram} (static constraint) applied under the same decoding settings.
This comparator does not use $\pi$ and does not intervene; it only blocks repeated 3-grams.
(Reported in Appendix~\ref{app:baseline}.)

\section{Results}
\subsection{Main results: collapse reduction across pillars}
Table~\ref{tab:main} summarizes outcomes across all five pillars ($n=500$ each).
Controlled decoding reduces collapse in every pillar with small runtime overhead.

\begin{table}[t]
\centering
\small
\setlength{\tabcolsep}{3pt}
\caption{Collapse reduction across five pillars under baseline vs.\ controlled decoding (Gemma 7B, Ollama, 512 tokens fixed).
Rates show Wilson 95\% CIs; $\Delta$/RR/OR use paired bootstrap over seeds (B=20{,}000).}
\label{tab:main}
\resizebox{\linewidth}{!}{%
\begin{tabular}{lrrrrrrrrrrrrr}
\toprule
Pillar & n & collapse\_base & collapse\_ctrl & $\Delta$ & $\Delta$ 95\%CI & RR & RR 95\%CI & OR & OR 95\%CI & intv\_rate & rt\_base & rt\_ctrl & $\Delta$rt \\
\midrule
HUM  & 500 & 0.142 & 0.000 & -0.142 & [-0.172,-0.112] & 0.007 & [0.0058,0.0088] & 0.006 & [0.0048,0.0079] & 0.142 & 19.6 & 20.9 & 1.2 \\
STEM & 500 & 0.302 & 0.004 & -0.298 & [-0.338,-0.258] & 0.013 & [0.0031,0.0345] & 0.009 & [0.0021,0.0245] & 0.302 & 19.1 & 20.7 & 1.6 \\
ETH  & 500 & 0.010 & 0.000 & -0.010 & [-0.020,-0.002] & 0.091 & [0.0476,0.3333] & 0.090 & [0.0467,0.3327] & 0.010 & 20.4 & 20.3 & -0.1 \\
TEMP & 500 & 0.030 & 0.000 & -0.030 & [-0.046,-0.016] & 0.032 & [0.0213,0.0588] & 0.031 & [0.0203,0.0579] & 0.030 & 19.9 & 20.0 & 0.1 \\
META & 500 & 0.106 & 0.006 & -0.100 & [-0.126,-0.074] & 0.057 & [0.0099,0.1277] & 0.051 & [0.0089,0.1170] & 0.106 & 19.8 & 20.5 & 0.7 \\
\bottomrule
\end{tabular}
}
\end{table}

\subsection{Cost and selectivity}
Interventions are selective: the fraction of runs with any intervention (\texttt{intv\_rate})
tracks baseline difficulty by pillar.
Runtime overhead is modest (Table~\ref{tab:main}).

\section{Discussion}
Our results indicate that collapse can be mitigated substantially with decode-time monitoring and minimal interventions,
without modifying model weights.
Limitations include dependence on the inference stack (Ollama) and fixed generation length (512 tokens).
Future work includes broader model coverage and longer contexts under the same paired-seed evaluation.

\section{Related Work}
Degeneration in neural text generation under sampling has been widely discussed, including repetition and incoherence under nucleus sampling.
Common decode-time mitigations include repetition penalties and $n$-gram blocking, but these are static rules and do not adapt online.
Other approaches add self-monitoring, critique/rewrite, or multi-pass decoding, typically increasing cost.
Our approach complements these by using lightweight online signals and minimal interventions within a single-pass generation loop,
evaluated under a paired-seed protocol that isolates the effect of intervention.

\section{Reproducibility Statement}
All results are generated via an auditable pipeline:
raw JSONL logs $\rightarrow$ run-level summaries (CSV) $\rightarrow$ Table~1 (CSV/MD).
No manual editing of results is performed.
Paired seeds ensure baseline and controlled runs share identical prompt selection.

% --- NeurIPS Paper Checklist ---
% IMPORTANT: Do not remove. Paste the official checklist block from the NeurIPS 2025 kit here.

\appendix

\section{A. Stress-probe (STEM, 30 seeds)}
\label{app:stress}
(Placeholder) Include a stress-probe subset and show detector activations before collapse.
Add 1--2 paired examples (baseline vs controlled) for the same seed/prompt.

\section{B. Collapse-reason breakdown (Table 2)}
\label{app:reasons}
(Placeholder) Auto-generate breakdown counts by collapse reason (REP vs TOO\_SHORT vs $\pi$-trajectory)
for each pillar and mode.

\section{C. JSONL schema summary}
(Placeholder) Document meta/step/final keys and invariants (1 run = 1 jsonl; no manual edits).

\section{D. Lightweight baseline comparator: no-repeat 3-gram}
\label{app:baseline}
(Placeholder) Report collapse rates under no-repeat 3-gram constraint to show that static blocking alone
does not match the online intervention gains across pillars.

\bibliographystyle{plain}
\end{document}
